\documentclass[10pt,aspectratio=169]{beamer}

% ------------------------------------------------------------------
% Presentació adaptada: Descriptiva bivariant amb orientació computacional
% Grau en Enginyeria Informàtica - Estadística
% ------------------------------------------------------------------

\usepackage[utf8]{inputenc}
\usepackage[T1]{fontenc}
\usepackage[catalan]{babel}
\usepackage{amsmath,amssymb}
\usepackage{booktabs}
\usepackage{array}
\usepackage{tikz}
\usepackage{pgfplots}
\usepackage{xcolor}
\usepackage{listings}
\pgfplotsset{compat=1.18}

% ------------------------------------------------------------------
% Estil visual coherent amb Descriptiva univariant computacional
% ------------------------------------------------------------------
\definecolor{UABGreen}{RGB}{0,103,71}
\definecolor{UABLight}{RGB}{224,241,236}
\definecolor{DeepBlue}{RGB}{20,65,110}
\definecolor{AccentOrange}{RGB}{230,126,34}
\definecolor{SoftGray}{RGB}{245,247,250}
\definecolor{CodeBack}{RGB}{248,248,248}
\definecolor{CodeFrame}{RGB}{210,215,220}

\usetheme{Madrid}
\usecolortheme[named=UABGreen]{structure}
\setbeamercolor{title}{fg=white,bg=UABGreen}
\setbeamercolor{frametitle}{fg=white,bg=UABGreen}
\setbeamercolor{block title}{fg=white,bg=DeepBlue}
\setbeamercolor{block body}{fg=black,bg=SoftGray}
\setbeamercolor{alerted text}{fg=AccentOrange}
\setbeamertemplate{navigation symbols}{}
\setbeamertemplate{itemize item}{\color{UABGreen}$\blacktriangleright$}
\setbeamertemplate{itemize subitem}{\color{AccentOrange}$\bullet$}
\setbeamertemplate{enumerate item}{\color{UABGreen}\insertenumlabel.}

\setbeamertemplate{footline}{%
  \leavevmode%
  \hbox{%
  \begin{beamercolorbox}[wd=.72\paperwidth,ht=2.4ex,dp=1ex,leftskip=1.4em]{author in head/foot}%
    Estadística -- Grau en Enginyeria Informàtica
  \end{beamercolorbox}%
  \begin{beamercolorbox}[wd=.28\paperwidth,ht=2.4ex,dp=1ex,center]{date in head/foot}%
    \insertframenumber{} / \inserttotalframenumber
  \end{beamercolorbox}}%
  \vskip0pt%
}

\lstdefinelanguage{R}{%
  morekeywords={if,else,repeat,while,function,for,in,next,break,TRUE,FALSE,NULL,NA,NaN,Inf,
    library,mean,median,sd,var,summary,table,prop.table,xtabs,addmargins,round,
    ggplot,aes,geom_point,geom_smooth,geom_col,geom_tile,facet_wrap,scale_fill_gradient,
    labs,theme_minimal,summarise,group_by,mutate,read_csv,filter,arrange,set.seed,
    rnorm,runif,cor,cov,lm,predict,residuals,broom,augment,glance,tibble,rep,seq},
  sensitive=true,
  morecomment=[l]\#,
  morestring=[b]",
  morestring=[b]'
}
\lstset{%
  language=R,
  basicstyle=\ttfamily\scriptsize,
  keywordstyle=\color{DeepBlue}\bfseries,
  commentstyle=\color{gray},
  stringstyle=\color{AccentOrange},
  backgroundcolor=\color{CodeBack},
  frame=single,
  rulecolor=\color{CodeFrame},
  breaklines=true,
  columns=fullflexible,
  keepspaces=true,
  showstringspaces=false,
  literate={à}{{\`a}}1 {è}{{\`e}}1 {é}{{\'e}}1 {í}{{\'i}}1 {ï}{{\"i}}1 {ò}{{\`o}}1 {ó}{{\'o}}1 {ú}{{\'u}}1 {ü}{{\"u}}1 {ç}{{\c{c}}}1 {À}{{\`A}}1 {È}{{\`E}}1 {É}{{\'E}}1 {Í}{{\'I}}1 {Ò}{{\`O}}1 {Ó}{{\'O}}1 {Ú}{{\'U}}1
}

\newcommand{\R}{\textsf{R}}
\newcommand{\RStudio}{\textsf{RStudio/Posit}}
\newcommand{\GitHub}{\textsf{GitHub}}
\newcommand{\important}[1]{\textcolor{UABGreen}{\textbf{#1}}}
\newcommand{\exemple}[1]{\textcolor{AccentOrange}{\textbf{#1}}}

% ------------------------------------------------------------------
% Dades de la presentació
% ------------------------------------------------------------------
\title[Descriptiva bivariant computacional]{Estadística descriptiva bivariant}
\author{Estadística -- Grau en Enginyeria Informàtica}
\date{Curs 2026--27}

\begin{document}

% ------------------------------------------------------------------
\begin{frame}[plain]
  \titlepage
  \vspace{-0.5cm}
  \begin{center}
  \small
  \textcolor{DeepBlue}{Dues variables, una pregunta: hi ha relació?}
  \end{center}
\end{frame}

% ------------------------------------------------------------------
\begin{frame}{Per què descriptiva bivariant en informàtica?}
\begin{columns}[T]
\column{0.55\textwidth}
\begin{itemize}
  \item Molts problemes computacionals no depenen d'una sola variable.
  \item Volem entendre \important{relacions} entre variables:
  \begin{itemize}
    \item navegador i estat d'una petició,
    \item sistema operatiu i tipus d'error,
    \item mida de l'entrada i temps d'execució,
    \item càrrega del servidor i latència,
    \item nombre d'usuaris i consum de memòria.
  \end{itemize}
  \item La descriptiva bivariant és el pas previ a la inferència, la regressió i el machine learning.
\end{itemize}
\column{0.4\textwidth}
\begin{block}{Idea clau}
Descriure dues variables és buscar patrons, dependències i anomalies en un sistema.
\end{block}
\vspace{0.2cm}
\begin{tikzpicture}
\node[draw=UABGreen,rounded corners,fill=UABLight,text width=4.5cm,align=center,inner sep=8pt] {
Variable A $+$ Variable B $\rightarrow$ Relació $\rightarrow$ Interpretació $\rightarrow$ Decisió
};
\end{tikzpicture}
\end{columns}
\end{frame}

% ------------------------------------------------------------------
\begin{frame}{Objectius del tema}
Al final d'aquest tema hauríeu de poder:
\begin{enumerate}
  \item Construir i interpretar taules de contingència.
  \item Calcular freqüències marginals, relatives i condicionades amb \R.
  \item Visualitzar relacions entre dues variables qualitatives.
  \item Construir i interpretar núvols de punts per a dues variables numèriques.
  \item Calcular covariància, correlació i recta de regressió.
  \item Interpretar prediccions i residus amb prudència.
  \item Documentar l'anàlisi en un informe reproduïble amb codi i resultats.
\end{enumerate}
\vspace{0.2cm}
\begin{block}{Enfocament computacional}
Les fórmules són importants, però també ho és saber convertir-les en codi fiable, gràfics clars i conclusions útils.
\end{block}
\end{frame}

% ------------------------------------------------------------------
\begin{frame}{Mapa del tema}
\begin{center}
\begin{tikzpicture}[node distance=1.35cm, every node/.style={font=\small}]
\node[draw,rounded corners,fill=UABLight,text width=3.2cm,align=center,minimum height=0.8cm] (qq) {Dues variables qualitatives};
\node[draw,rounded corners,fill=UABLight,right of=qq,xshift=3.0cm,text width=3.2cm,align=center,minimum height=0.8cm] (nn) {Dues variables numèriques};
\node[draw,rounded corners,fill=SoftGray,below of=qq,yshift=-0.2cm,text width=3.2cm,align=center,minimum height=0.8cm] (taules) {Taules de contingència};
\node[draw,rounded corners,fill=SoftGray,below of=nn,yshift=-0.2cm,text width=3.2cm,align=center,minimum height=0.8cm] (scatter) {Núvols de punts};
\node[draw,rounded corners,fill=SoftGray,below of=taules,yshift=-0.2cm,text width=3.2cm,align=center,minimum height=0.8cm] (cond) {Freqüències condicionades};
\node[draw,rounded corners,fill=SoftGray,below of=scatter,yshift=-0.2cm,text width=3.2cm,align=center,minimum height=0.8cm] (corr) {Correlació i regressió};
\draw[->,thick,UABGreen] (qq) -- (taules);
\draw[->,thick,UABGreen] (taules) -- (cond);
\draw[->,thick,UABGreen] (nn) -- (scatter);
\draw[->,thick,UABGreen] (scatter) -- (corr);
\draw[->,thick,AccentOrange] (cond) -- ++(0,-0.7) -| (corr);
\end{tikzpicture}
\end{center}
\vspace{0.25cm}
\begin{alertblock}{Pregunta transversal}
Quan veiem una associació, és real, és rellevant, és lineal, és causal o és un artefacte de les dades?
\end{alertblock}
\end{frame}

% ------------------------------------------------------------------
\section{Dues variables qualitatives}

\begin{frame}{Dues variables qualitatives: taula de contingència}
Considerem dues variables qualitatives:
\begin{itemize}
  \item $A$ amb modalitats $A_1,A_2,\ldots,A_k$.
  \item $B$ amb modalitats $B_1,B_2,\ldots,B_\ell$.
\end{itemize}
\vspace{0.2cm}
\begin{columns}[T]
\column{0.52\textwidth}
\centering
\small
\begin{tabular}{lcccc}
\toprule
 & $B_1$ & $B_2$ & $\cdots$ & $B_\ell$ \\
\midrule
$A_1$ & $n_{11}$ & $n_{12}$ & $\cdots$ & $n_{1\ell}$ \\
$A_2$ & $n_{21}$ & $n_{22}$ & $\cdots$ & $n_{2\ell}$ \\
$\vdots$ & $\vdots$ & $\vdots$ & $\ddots$ & $\vdots$ \\
$A_k$ & $n_{k1}$ & $n_{k2}$ & $\cdots$ & $n_{k\ell}$ \\
\bottomrule
\end{tabular}
\column{0.43\textwidth}
\begin{block}{Interpretació}
$n_{ij}$ és el nombre d'observacions que tenen simultàniament la modalitat $A_i$ i la modalitat $B_j$.
\end{block}
\begin{exampleblock}{Exemple informàtic}
$A =$ navegador, $B =$ estat de la petició: \texttt{OK}, \texttt{timeout}, \texttt{error 500}.
\end{exampleblock}
\end{columns}
\end{frame}

% ------------------------------------------------------------------
\begin{frame}{Exemple: peticions web per navegador i estat}
\begin{columns}[T]
\column{0.50\textwidth}
\small
\begin{tabular}{lrrrr}
\toprule
Navegador & OK & Timeout & Error & Total \\
\midrule
Chrome  & 820 & 25 & 35 & 880 \\
Firefox & 430 & 18 & 22 & 470 \\
Safari  & 260 & 12 & 28 & 300 \\
Edge    & 210 & 15 & 25 & 250 \\
\midrule
Total   & 1720 & 70 & 110 & 1900 \\
\bottomrule
\end{tabular}
\vspace{0.3cm}
\begin{block}{Lectura ràpida}
La major part de peticions són correctes, però la proporció d'errors no és igual en tots els navegadors.
\end{block}
\column{0.45\textwidth}
\centering
\begin{tikzpicture}
\begin{axis}[
  ybar stacked,
  width=6.0cm,
  height=4.3cm,
  ymin=0,
  ylabel={Peticions},
  symbolic x coords={Chrome,Firefox,Safari,Edge},
  xtick=data,
  x tick label style={rotate=25,anchor=east,font=\scriptsize},
  legend style={font=\scriptsize,at={(0.5,-0.25)},anchor=north,legend columns=3},
  bar width=12pt,
]
\addplot+[fill=UABGreen!65,draw=UABGreen] coordinates {(Chrome,820) (Firefox,430) (Safari,260) (Edge,210)};
\addplot+[fill=AccentOrange!65,draw=AccentOrange] coordinates {(Chrome,25) (Firefox,18) (Safari,12) (Edge,15)};
\addplot+[fill=DeepBlue!55,draw=DeepBlue] coordinates {(Chrome,35) (Firefox,22) (Safari,28) (Edge,25)};
\legend{OK,Timeout,Error}
\end{axis}
\end{tikzpicture}
\end{columns}
\end{frame}

% ------------------------------------------------------------------
\begin{frame}{Freqüències marginals}
Les freqüències marginals són els totals de files i columnes.
\vspace{0.2cm}
\begin{columns}[T]
\column{0.48\textwidth}
\begin{block}{Marge dret}
\[
 n_{i\cdot}=\sum_{j=1}^{\ell} n_{ij}
\]
Freqüència total de la modalitat $A_i$.
\end{block}
\column{0.48\textwidth}
\begin{block}{Marge inferior}
\[
 n_{\cdot j}=\sum_{i=1}^{k} n_{ij}
\]
Freqüència total de la modalitat $B_j$.
\end{block}
\end{columns}
\vspace{0.3cm}
\begin{exampleblock}{En l'exemple}
Hi ha 880 peticions de Chrome i 110 peticions amb error. La grandària total és $n=1900$.
\end{exampleblock}
\end{frame}

% ------------------------------------------------------------------
\begin{frame}{Freqüències relatives}
Les freqüències relatives s'obtenen dividint per la grandària total $n$:
\[
 f_{ij}=f(A_i\cap B_j)=\frac{n_{ij}}{n},\qquad
 f_{i\cdot}=\frac{n_{i\cdot}}{n},\qquad
 f_{\cdot j}=\frac{n_{\cdot j}}{n}.
\]
\vspace{0.15cm}
\begin{columns}[T]
\column{0.48\textwidth}
\begin{block}{Per què són útils?}
\begin{itemize}
  \item Permeten comparar mostres de mida diferent.
  \item Es poden llegir com a percentatges.
  \item Són la base de les probabilitats empíriques.
\end{itemize}
\end{block}
\column{0.48\textwidth}
\begin{exampleblock}{Exemple}
\[
 f(\text{Chrome i Error})=\frac{35}{1900}=0.0184.
\]
Aproximadament un 1.84\% de totes les peticions són errors amb Chrome.
\end{exampleblock}
\end{columns}
\end{frame}

% ------------------------------------------------------------------
\begin{frame}{Freqüències condicionades}
Quan volem comparar correctament, sovint hem de condicionar.
\vspace{0.15cm}
\begin{columns}[T]
\column{0.48\textwidth}
\begin{block}{Condicionar per columnes}
\[
 f(A_i\mid B_j)=\frac{n_{ij}}{n_{\cdot j}}
\]
Distribució de $A$ dins la subpoblació $B_j$.
\end{block}
\column{0.48\textwidth}
\begin{block}{Condicionar per files}
\[
 f(B_j\mid A_i)=\frac{n_{ij}}{n_{i\cdot}}
\]
Distribució de $B$ dins la subpoblació $A_i$.
\end{block}
\end{columns}
\vspace{0.25cm}
\begin{alertblock}{Missatge important}
No és el mateix $f(A_i\cap B_j)$ que $f(B_j\mid A_i)$. Una pregunta sobre risc, error o rendiment acostuma a ser una pregunta condicionada.
\end{alertblock}
\end{frame}

% ------------------------------------------------------------------
\begin{frame}{Exemple: percentatge d'error per navegador}
\begin{columns}[T]
\column{0.50\textwidth}
\small
\begin{tabular}{lrrr}
\toprule
Navegador & Total & Errors & $f(\text{Error}\mid\text{Navegador})$ \\
\midrule
Chrome  & 880 & 35 & 3.98\% \\
Firefox & 470 & 22 & 4.68\% \\
Safari  & 300 & 28 & 9.33\% \\
Edge    & 250 & 25 & 10.00\% \\
\bottomrule
\end{tabular}
\vspace{0.35cm}
\begin{block}{Lectura}
La freqüència absoluta d'errors és més gran a Chrome, però el percentatge d'error és més alt a Safari i Edge.
\end{block}
\column{0.45\textwidth}
\centering
\begin{tikzpicture}
\begin{axis}[
  ybar,
  width=6cm,
  height=4.2cm,
  ymin=0,ymax=12,
  ylabel={Error (\%)},
  symbolic x coords={Chrome,Firefox,Safari,Edge},
  xtick=data,
  x tick label style={rotate=25,anchor=east,font=\scriptsize},
  bar width=14pt,
  nodes near coords,
  nodes near coords style={font=\scriptsize},
]
\addplot+[fill=AccentOrange!70,draw=AccentOrange] coordinates {(Chrome,3.98) (Firefox,4.68) (Safari,9.33) (Edge,10.00)};
\end{axis}
\end{tikzpicture}
\end{columns}
\end{frame}

% ------------------------------------------------------------------
\begin{frame}[fragile]{Taules de contingència amb R}
\begin{lstlisting}
# dades: una fila per peticio
logs <- read_csv("data/logs_web.csv")

# taula absoluta
tab <- table(logs$navegador, logs$estat)
addmargins(tab)

# taula relativa global
round(prop.table(tab), 3)

# percentatges per fila: estat condicionat al navegador
round(prop.table(tab, margin = 1), 3)

# percentatges per columna: navegador condicionat a l'estat
round(prop.table(tab, margin = 2), 3)
\end{lstlisting}
\vspace{0.2cm}
\begin{block}{Objectiu pràctic}
Que la mateixa taula es pugui regenerar quan arriben nous logs o quan canviem el període d'anàlisi.
\end{block}
\end{frame}

% ------------------------------------------------------------------
\begin{frame}[fragile]{Visualitzar una taula: del número al patró}
\begin{columns}[T]
\column{0.52\textwidth}
\begin{lstlisting}
library(ggplot2)

logs |>
  count(navegador, estat) |>
  group_by(navegador) |>
  mutate(prop = n / sum(n)) |>
  ggplot(aes(navegador, estat, fill = prop)) +
  geom_tile() +
  labs(fill = "Freq. cond.",
       x = "Navegador", y = "Estat") +
  theme_minimal()
\end{lstlisting}
\column{0.43\textwidth}
\centering
\begin{tikzpicture}
\foreach \x/\lab in {0/Chrome,1.2/Firefox,2.4/Safari,3.6/Edge} {
  \node[font=\scriptsize,rotate=35,anchor=east] at (\x+0.35,-0.2) {\lab};
}
\foreach \y/\lab in {0/OK,0.8/Timeout,1.6/Error} {
  \node[font=\scriptsize,anchor=east] at (-0.15,\y+0.35) {\lab};
}
\foreach \x/\y/\c in {0/0/UABGreen!80,1.2/0/UABGreen!75,2.4/0/UABGreen!55,3.6/0/UABGreen!50,
0/0.8/AccentOrange!20,1.2/0.8/AccentOrange!25,2.4/0.8/AccentOrange!28,3.6/0.8/AccentOrange!35,
0/1.6/DeepBlue!25,1.2/1.6/DeepBlue!30,2.4/1.6/DeepBlue!50,3.6/1.6/DeepBlue!55} {
  \draw[fill=\c,draw=white] (\x,\y) rectangle +(0.9,0.7);
}
\node[align=center,font=\small,text width=4.5cm] at (1.9,2.7) {Mapa de calor de freqüències condicionades};
\end{tikzpicture}
\end{columns}
\end{frame}

% ------------------------------------------------------------------
\begin{frame}{Associació entre dues variables qualitatives}
\begin{itemize}
  \item Si la distribució de $B$ és semblant per a tots els nivells de $A$, direm que no hi ha un patró clar d'associació.
  \item Si les freqüències condicionades canvien molt segons $A$, hi ha evidència descriptiva d'associació.
  \item La inferència estadística permetrà decidir si aquesta diferència pot ser compatible amb l'atzar.
\end{itemize}
\vspace{0.3cm}
\begin{exampleblock}{Pregunta computacional}
El tipus d'error depèn del navegador, del sistema operatiu o de la versió desplegada?
\end{exampleblock}
\begin{alertblock}{Prudència}
Una taula de contingència pot mostrar associació, però no demostra causalitat.
\end{alertblock}
\end{frame}

% ------------------------------------------------------------------
\section{Dues variables numèriques}

\begin{frame}{Dues variables numèriques: preguntes bàsiques}
En l'estudi descriptiu de dues variables numèriques volem respondre:
\begin{itemize}
  \item Tenen aquestes variables alguna relació?
  \item La relació és positiva, negativa o nul·la?
  \item És aproximadament lineal o no lineal?
  \item Hi ha observacions anòmales?
  \item Podem predir una variable a partir de l'altra?
\end{itemize}
\vspace{0.25cm}
\begin{exampleblock}{Exemples informàtics}
\begin{itemize}
  \item Mida de l'entrada i temps d'execució.
  \item Nombre d'usuaris concurrents i latència.
  \item Mida d'un fitxer i temps de compressió.
  \item Nombre de tests i temps de CI/CD.
\end{itemize}
\end{exampleblock}
\end{frame}

% ------------------------------------------------------------------
\begin{frame}{La dada bivariant: parelles $(x_i,y_i)$}
Per a cada observació $i$ tenim dues mesures:
\[
 (x_i,y_i),\qquad i=1,\ldots,n.
\]
\vspace{0.2cm}
\begin{columns}[T]
\column{0.48\textwidth}
\small
\begin{tabular}{rrrr}
\toprule
Execució & mida $x$ & temps $y$ & memòria \\
\midrule
1 & 1000 & 12.4 & 45 \\
2 & 2000 & 19.8 & 51 \\
3 & 3000 & 29.1 & 59 \\
4 & 4000 & 38.7 & 67 \\
5 & 5000 & 51.0 & 77 \\
\bottomrule
\end{tabular}
\column{0.48\textwidth}
\begin{block}{Visualització natural}
Cada parella $(x_i,y_i)$ es representa com un punt en el pla.
\end{block}
\begin{tikzpicture}
\begin{axis}[width=5.2cm,height=3.5cm,xlabel={mida},ylabel={temps},grid=major]
\addplot+[only marks,mark=*,color=UABGreen] coordinates {(1,12.4) (2,19.8) (3,29.1) (4,38.7) (5,51)};
\end{axis}
\end{tikzpicture}
\end{columns}
\end{frame}

% ------------------------------------------------------------------
\begin{frame}{Núvol de punts}
\begin{columns}[T]
\column{0.52\textwidth}
El \important{núvol de punts} és la representació gràfica de totes les parelles $(x_i,y_i)$.
\vspace{0.2cm}
\begin{itemize}
  \item Mostra tendències globals.
  \item Detecta relacions no lineals.
  \item Identifica valors atípics.
  \item Evita confiar només en una única mesura numèrica.
\end{itemize}
\column{0.43\textwidth}
\centering
\begin{tikzpicture}
\begin{axis}[width=5.7cm,height=4.0cm,xlabel={mida entrada},ylabel={temps ms},grid=major]
\addplot+[only marks,mark=*,color=UABGreen] coordinates {
(1,9) (2,14) (3,18) (4,22) (5,28) (6,31) (7,36) (8,42) (9,46) (10,52)
(2,17) (4,24) (6,34) (8,39) (10,55)};
\end{axis}
\end{tikzpicture}
\end{columns}
\end{frame}

% ------------------------------------------------------------------
\begin{frame}{Patrons habituals en un núvol de punts}
\begin{columns}[T]
\column{0.5\textwidth}
\centering
\begin{tikzpicture}
\begin{axis}[width=5.1cm,height=3.4cm,title={Relació positiva},xtick=\empty,ytick=\empty]
\addplot+[only marks,color=UABGreen] coordinates {(1,2) (2,3.1) (3,3.6) (4,5.2) (5,5.7) (6,6.5) (7,8)};
\end{axis}
\end{tikzpicture}
\vspace{0.1cm}
\begin{tikzpicture}
\begin{axis}[width=5.1cm,height=3.4cm,title={Sense relació clara},xtick=\empty,ytick=\empty]
\addplot+[only marks,color=DeepBlue] coordinates {(1,5) (2,2) (3,7) (4,3.2) (5,8) (6,4) (7,6)};
\end{axis}
\end{tikzpicture}
\column{0.5\textwidth}
\centering
\begin{tikzpicture}
\begin{axis}[width=5.1cm,height=3.4cm,title={Relació negativa},xtick=\empty,ytick=\empty]
\addplot+[only marks,color=AccentOrange] coordinates {(1,8) (2,6.8) (3,6) (4,4.8) (5,3.5) (6,3) (7,2)};
\end{axis}
\end{tikzpicture}
\vspace{0.1cm}
\begin{tikzpicture}
\begin{axis}[width=5.1cm,height=3.4cm,title={Relació no lineal},xtick=\empty,ytick=\empty]
\addplot+[only marks,color=UABGreen] coordinates {(1,7) (2,4) (3,2.5) (4,2) (5,2.5) (6,4) (7,7)};
\end{axis}
\end{tikzpicture}
\end{columns}
\end{frame}

% ------------------------------------------------------------------
\begin{frame}{Covariància}
\begin{columns}[T]
\column{0.52\textwidth}
Les mitjanes de les dues variables són:
\[
 \bar{x}=\frac{1}{n}\sum_{i=1}^n x_i,\qquad
 \bar{y}=\frac{1}{n}\sum_{i=1}^n y_i.
\]
La covariància mostral és:
\[
 S_{XY}=\frac{1}{n-1}\sum_{i=1}^n (x_i-\bar{x})(y_i-\bar{y}).
\]
\column{0.43\textwidth}
\begin{block}{Interpretació}
\begin{itemize}
  \item $S_{XY}>0$: tendència conjunta positiva.
  \item $S_{XY}<0$: tendència conjunta negativa.
  \item $S_{XY}\approx 0$: no hi ha tendència lineal clara.
\end{itemize}
\end{block}
\begin{alertblock}{Limitació}
La covariància depèn de les unitats de mesura i és difícil de comparar entre problemes.
\end{alertblock}
\end{columns}
\end{frame}

% ------------------------------------------------------------------
\begin{frame}{Correlació de Pearson}
La correlació de Pearson normalitza la covariància:
\[
 r_{XY}=\frac{S_{XY}}{S_XS_Y},\qquad -1\leq r_{XY}\leq 1.
\]
\vspace{0.2cm}
\begin{columns}[T]
\column{0.48\textwidth}
\begin{block}{Lectura del signe}
\begin{itemize}
  \item $r>0$: relació lineal positiva.
  \item $r<0$: relació lineal negativa.
  \item $r\approx 0$: no hi ha relació lineal clara.
\end{itemize}
\end{block}
\column{0.48\textwidth}
\begin{block}{Lectura de la magnitud}
\begin{itemize}
  \item $|r|$ proper a 1: relació lineal forta.
  \item $|r|$ proper a 0: relació lineal feble.
  \item $r^2$: proporció descriptiva de variabilitat explicada per la recta.
\end{itemize}
\end{block}
\end{columns}
\end{frame}

% ------------------------------------------------------------------
\begin{frame}[fragile]{Correlació amb R}
\begin{columns}[T]
\column{0.55\textwidth}
\begin{lstlisting}
bench <- read_csv("data/benchmark.csv")

# covariancia i correlacio
cov(bench$mida, bench$temps_ms)
cor(bench$mida, bench$temps_ms)

# grafic bivariant
ggplot(bench, aes(x = mida, y = temps_ms)) +
  geom_point() +
  geom_smooth(method = "lm", se = FALSE) +
  theme_minimal()
\end{lstlisting}
\column{0.4\textwidth}
\begin{block}{Bones pràctiques}
\begin{itemize}
  \item Mirar sempre el gràfic abans d'interpretar $r$.
  \item Revisar valors atípics.
  \item No confondre correlació amb causalitat.
  \item Informar unitats i context.
\end{itemize}
\end{block}
\end{columns}
\end{frame}

% ------------------------------------------------------------------
\begin{frame}{Correlació: advertiments importants}
\begin{columns}[T]
\column{0.48\textwidth}
\begin{alertblock}{1. Correlació no implica causalitat}
Que dues variables es moguin juntes no vol dir que una causi l'altra.
\end{alertblock}
\begin{alertblock}{2. Pearson mesura relació lineal}
Una relació no lineal pot tenir $r\approx 0$.
\end{alertblock}
\column{0.48\textwidth}
\begin{alertblock}{3. Els outliers poden dominar}
Un sol valor anòmal pot canviar molt la correlació.
\end{alertblock}
\begin{alertblock}{4. Barrejar grups pot enganyar}
La relació global pot ser diferent de la relació dins de cada grup.
\end{alertblock}
\end{columns}
\vspace{0.2cm}
\begin{exampleblock}{Exemple informàtic}
La relació entre mida d'entrada i temps pot dependre de l'algoritme, del maquinari i de la memòria cau.
\end{exampleblock}
\end{frame}

% ------------------------------------------------------------------
\section{Recta de regressió}

\begin{frame}{Recta de regressió: primera idea}
La recta de regressió resumeix una relació lineal entre dues variables:
\[
 \widehat{y}=a+bx.
\]
\vspace{0.1cm}
\begin{columns}[T]
\column{0.48\textwidth}
\begin{itemize}
  \item $x$: variable explicativa o predictora.
  \item $y$: variable resposta.
  \item $b$: canvi mitjà en $y$ per una unitat més de $x$.
  \item $a$: valor predit quan $x=0$.
\end{itemize}
\column{0.48\textwidth}
\begin{block}{Criteri}
Triem la recta que minimitza la suma de quadrats dels residus:
\[
 \sum_{i=1}^n \left[y_i-(a+bx_i)\right]^2.
\]
\end{block}
\end{columns}
\end{frame}

% ------------------------------------------------------------------
\begin{frame}{Coeficients de la recta}
La recta ajustada és:
\[
 \widehat{y}=a+bx.
\]
Els coeficients es calculen com:
\[
 b=\frac{S_{XY}}{S_X^2},\qquad a=\bar{y}-b\bar{x}.
\]
\vspace{0.2cm}
\begin{block}{Interpretació computacional}
Si $x$ és la mida de l'entrada i $y$ és el temps d'execució, $b$ estima l'increment mitjà de temps quan augmentem una unitat la mida d'entrada.
\end{block}
\begin{alertblock}{Atenció}
La recta és un resum descriptiu. No garanteix que el fenomen sigui realment lineal ni que pugui extrapolar-se fora del rang observat.
\end{alertblock}
\end{frame}

% ------------------------------------------------------------------
\begin{frame}{Exemple: mida de l'entrada i temps d'execució}
\begin{columns}[T]
\column{0.48\textwidth}
\small
\begin{tabular}{rrrr}
\toprule
$mida$ & $temps$ & $\widehat{temps}$ & residu \\
\midrule
1000 & 12.4 & 11.7 & 0.7 \\
2000 & 19.8 & 21.3 & -1.5 \\
3000 & 29.1 & 30.9 & -1.8 \\
4000 & 38.7 & 40.5 & -1.8 \\
5000 & 51.0 & 50.1 & 0.9 \\
6000 & 61.8 & 59.7 & 2.1 \\
\bottomrule
\end{tabular}
\vspace{0.25cm}
\begin{block}{Model descriptiu}
\[
 \widehat{temps}=2.1+0.0096\cdot mida.
\]
\end{block}
\column{0.48\textwidth}
\centering
\begin{tikzpicture}
\begin{axis}[width=6.1cm,height=4.2cm,xlabel={mida entrada},ylabel={temps ms},grid=major]
\addplot+[only marks,mark=*,color=UABGreen] coordinates {(1000,12.4) (2000,19.8) (3000,29.1) (4000,38.7) (5000,51) (6000,61.8)};
\addplot+[domain=900:6200,samples=2,color=AccentOrange,thick] {2.1 + 0.0096*x};
\end{axis}
\end{tikzpicture}
\end{columns}
\end{frame}

% ------------------------------------------------------------------
\begin{frame}[fragile]{Regressió amb R}
\begin{lstlisting}
bench <- read_csv("data/benchmark.csv")

model <- lm(temps_ms ~ mida, data = bench)
summary(model)

# coeficients de la recta
coef(model)

# prediccions i residus
bench$prediccio <- predict(model)
bench$residu <- residuals(model)

# grafic amb recta ajustada
ggplot(bench, aes(mida, temps_ms)) +
  geom_point() +
  geom_smooth(method = "lm", se = FALSE) +
  theme_minimal()
\end{lstlisting}
\end{frame}

% ------------------------------------------------------------------
\begin{frame}{Predicció i residus}
Per a un valor $x^*$, la predicció és:
\[
 \widehat{y}^*=a+bx^*.
\]
El residu d'una observació és:
\[
 e_i=y_i-\widehat{y}_i.
\]
\vspace{0.15cm}
\begin{columns}[T]
\column{0.48\textwidth}
\begin{block}{Quan confiar més en una predicció?}
\begin{itemize}
  \item $x^*$ és dins el rang observat.
  \item El núvol de punts és aproximadament lineal.
  \item No hi ha outliers dominants.
  \item Els residus no mostren patrons evidents.
\end{itemize}
\end{block}
\column{0.48\textwidth}
\begin{alertblock}{Quan desconfiar?}
\begin{itemize}
  \item Extrapolem fora del rang.
  \item La relació és clarament corba.
  \item Hi ha canvis de règim.
  \item Barrejar grups amaga diferències.
\end{itemize}
\end{alertblock}
\end{columns}
\end{frame}

% ------------------------------------------------------------------
\begin{frame}{Mirar residus: detectar el que la recta no explica}
\begin{columns}[T]
\column{0.48\textwidth}
\centering
\begin{tikzpicture}
\begin{axis}[width=5.6cm,height=3.8cm,xlabel={mida},ylabel={residu},grid=major]
\addplot+[only marks,color=DeepBlue] coordinates {(1,0.7) (2,-1.5) (3,-1.8) (4,-1.8) (5,0.9) (6,2.1)};
\addplot+[domain=0.5:6.5,samples=2,color=AccentOrange,thick] {0};
\end{axis}
\end{tikzpicture}
\column{0.48\textwidth}
\begin{block}{Interpretació}
\begin{itemize}
  \item Residus al voltant de zero: ajust raonable.
  \item Patró corbat: la recta és massa simple.
  \item Residus molt grans: possibles outliers.
  \item Variabilitat creixent: la incertesa depèn de $x$.
\end{itemize}
\end{block}
\end{columns}
\end{frame}

% ------------------------------------------------------------------
\begin{frame}{De la descriptiva bivariant al modelatge}
\begin{center}
\begin{tikzpicture}[node distance=1.55cm, every node/.style={font=\small}]
\node[draw,rounded corners,fill=UABLight,minimum width=2.4cm,minimum height=0.8cm] (plot) {Gràfic};
\node[draw,rounded corners,fill=UABLight,right of=plot,xshift=1.6cm,minimum width=2.4cm,minimum height=0.8cm] (summary) {Mesura};
\node[draw,rounded corners,fill=UABLight,right of=summary,xshift=1.6cm,minimum width=2.4cm,minimum height=0.8cm] (model) {Model};
\node[draw,rounded corners,fill=UABLight,right of=model,xshift=1.6cm,minimum width=2.4cm,minimum height=0.8cm] (decision) {Decisió};
\draw[->,thick,UABGreen] (plot) -- (summary);
\draw[->,thick,UABGreen] (summary) -- (model);
\draw[->,thick,UABGreen] (model) -- (decision);
\end{tikzpicture}
\end{center}
\vspace{0.3cm}
\begin{columns}[T]
\column{0.48\textwidth}
\begin{block}{Qualitatives}
Taules, percentatges condicionats i comparació de perfils.
\end{block}
\column{0.48\textwidth}
\begin{block}{Numèriques}
Núvols de punts, correlació, regressió i residus.
\end{block}
\end{columns}
\vspace{0.25cm}
\begin{exampleblock}{Pont amb temes futurs}
Inferència, contrasts d'hipòtesis, regressió, classificació i avaluació de models.
\end{exampleblock}
\end{frame}

\end{document}
